% AWM Abstract — CoRL 2026
% Last updated: 2026-04-02

\begin{abstract}
Morphologically adaptive robots promise the versatility of legged locomotion and the
efficiency of wheeled movement, but typically require depth cameras and complex algorithms
to determine when to reconfigure. We argue that robot design choices directly determine
algorithmic complexity — a robot built with learnability in mind can reduce a hard
perception-and-control problem to a simple learned policy. We present AWM (All-Wheel-Morph),
a low-cost open-source wheeled-legged robot where each wheel extends into a leg, designed
to make morphological adaptation learnable: actuators with built-in torque observability,
a continuous morphology range, and sim-to-real transferable dynamics. Leveraging torque
feedback natively available from AWM's actuators, we propose \textbf{Active Morphological
Probing} --- an LSTM policy that physically senses terrain by extending legs and reading
contact torque, adapting morphology without any depth camera. Without exteroception, our
policy matches camera-equipped baselines in terrain traversal while a fixed-morphology
wheels-only policy plateaus at half the curriculum difficulty regardless of training
duration, confirming that adaptive morphology is the key enabler. Our approach requires
no teacher-student distillation, no privileged training information, and deploys as a
single network. We fully open-source AWM --- including hardware design, bill of materials,
simulation assets, training code, and trained policies --- so that anyone can build, train,
and deploy a morphologically adaptive robot with off-the-shelf components under \$2,500.
\end{abstract}
